%% Chronogramme des fonctions ET et OU pour les quatre combinaisons de deux entrées.
%% Le ET n'est haut que si les deux entrées sont hautes ; le OU dès que l'une l'est.
\documentclass[tikz,border=4pt]{standalone}
\usepackage{liaisons}
\begin{document}
\begin{tikzpicture}[x=1cm,y=1cm,
  sig/.style={line width=1.4pt, line join=round},
  ax/.style={black!35, line width=0.5pt, -{Stealth[length=4pt]}},
  sep/.style={black!45, dash pattern=on 2pt off 2pt, line width=0.5pt},
  nom/.style={font=\small, anchor=east, inner sep=4pt},
  et/.style={font=\scriptsize, anchor=south, inner sep=1pt}]

%% ================== axes de temps et séparateurs ======================
\foreach \y in {0,1,2,3} { \draw[ax] (-0.1,\y) -- (5.15,\y); }
\foreach \x in {0,1.2,2.4,3.6,4.8} { \draw[sep] (\x,-0.25) -- (\x,3.95); }

%% ================== entrées ===========================================
\draw[sig, solideB] (0,3) -- (2.4,3) -- (2.4,3.5) -- (4.8,3.5);
\draw[sig, solideB] (0,2) -- (1.2,2) -- (1.2,2.5) -- (2.4,2.5) -- (2.4,2)
                    -- (3.6,2) -- (3.6,2.5) -- (4.8,2.5);
%% ================== sorties ===========================================
\draw[sig, solideA] (0,1) -- (3.6,1) -- (3.6,1.5) -- (4.8,1.5);
\draw[sig, solideA] (0,0) -- (1.2,0) -- (1.2,0.5) -- (4.8,0.5);

%% ================== états logiques ====================================
\foreach \y/\liste in {3/{0,0,1,1}, 2/{0,1,0,1}, 1/{0,0,0,1}, 0/{0,1,1,1}} {
  \foreach \v [count=\i from 0] in \liste {
    \node[et] at ({(\i+0.5)*1.2},{\y+0.5*\v+0.06}) {\v}; } }

%% ================== noms des signaux ==================================
\node[nom] at (-0.1,3.25) {$e_1$};
\node[nom] at (-0.1,2.25) {$e_2$};
\node[nom, text=solideA] at (-0.1,1.25) {ET};
\node[nom, text=solideA] at (-0.1,0.25) {OU};
\node[font=\small, anchor=north east, inner sep=2pt] at (5.15,-0.05) {$t$};
\end{tikzpicture}
\end{document}
